\section{The homogeneous-space realization}

Sections~3 and~4 constructed the graph

\[
X=\KC(L(P))
\]

as the canonical bipartite double cover of the line graph of the Petersen
graph.

In this section we temporarily set that construction aside.

Instead we begin with a transitive permutation action of the symmetric
group \(S_{5}\) on a homogeneous space of degree thirty and study the
orbital graphs arising from that action.

Only at the end of the section do we identify the resulting graph with the
graph constructed previously.

\subsection{The degree-thirty coset actions}

The symmetric group

\[
S_{5}
\]

contains two conjugacy classes of Klein four subgroups, introduced in
Section~2.6.

These determine the homogeneous spaces

\[
\Omega_{\mathrm{even}}
=
S_{5}/V_{4,\mathrm{even}}
\]

and

\[
\Omega_{\mathrm{mixed}}
=
S_{5}/V_{4,\mathrm{mixed}}.
\]

Each action has degree

\[
30,
\]

since both Klein four subgroups have index thirty.

Although these actions have the same degree, they are not permutation
isomorphic because the corresponding subgroups are not conjugate in
\(S_{5}\).

Consequently they define two distinct homogeneous spaces equipped with
transitive actions of the symmetric group.

% STATUS: Conceptual


\subsection{The \(A_{5}\)-orbit distinction}

We compare the restrictions of the two degree-thirty actions to the
alternating subgroup \(A_{5}\). Their orbit structures are different.

\begin{proposition}

The restriction of the action on

\[
\Omega_{\mathrm{even}}
\]

to \(A_{5}\) has exactly two orbits, each of size fifteen.

The restriction of the action on

\[
\Omega_{\mathrm{mixed}}
\]

remains transitive.

\end{proposition}

\begin{proof}

The subgroup

\[
V_{4,\mathrm{even}}
\]

is contained in

\[
A_{5},
\]

so

\[
A_{5}
\]

acts on the coset space

\[
S_{5}/V_{4,\mathrm{even}}
\]

with two orbits corresponding to the two cosets of

\[
A_{5}
\]

in

\[
S_{5}.
\]

Each orbit therefore has size

\[
|A_{5}:V_{4,\mathrm{even}}|
=
60/4
=
15.
\]

In contrast,

\[
V_{4,\mathrm{mixed}}
\]

contains odd permutations and is not contained in

\[
A_{5}.
\]

Consequently,

\[
A_{5}V_{4,\mathrm{mixed}}
=
S_{5},
\]

and the restriction of the action remains transitive.

\end{proof}

Thus the two degree-thirty homogeneous spaces are already distinguished at
the level of the alternating subgroup.

The homogeneous space

\[
\Omega_{\mathrm{even}}
\]

naturally decomposes into two equal halves,

whereas

\[
\Omega_{\mathrm{mixed}}
\]

does not.

% STATUS: Conceptual


\subsection{Certified quartic orbital census}

% STATUS: Certified computation
% SOURCE: orbital_census_certificate
% Theorem established by exhaustive finite computation.

The action of

\[
S_{5}
\]

on

\[
\Omega_{\mathrm{even}}
=
S_{5}/V_{4,\mathrm{even}}
\]

admits a finite family of self-paired orbital graphs.

The complete quartic orbital census was performed by exhaustive
computation and is recorded in the certified artifacts described in
Section~9.

\begin{theorem}[Certified quartic orbital census]
\label{thm:quartic-orbital-census}

Among the self-paired orbital graphs arising from the action of

\[
S_{5}
\]

on

\[
S_{5}/V_{4,\mathrm{even}},
\]

there exist exactly four quartic orbital graphs.

Three are disconnected.

Exactly one is connected.

\end{theorem}

The four orbital graphs are summarized in
Table~\ref{tab:quartic-orbitals}.

\begin{table}[h]
\centering
\begin{tabular}{llll}
\toprule
Components &
Triangle-free &
Bipartite &
Status\\
\midrule
15+15 & No & No & Certified\\
30 & Yes & Yes & Certified\\
15+15 & No & No & Certified\\
15+15 & No & No & Certified\\
\bottomrule
\end{tabular}
\caption{Certified quartic self-paired orbital census for
\(S_{5}/V_{4,\mathrm{even}}\).}
\label{tab:quartic-orbitals}
\end{table}

No conceptual proof of this classification is presently known.
The theorem records the exact outcome of the certified finite census.


\subsection{The connected triangle-free orbital}

% STATUS: Conceptual deduction from certified computation

The graph constructed in Section~3 is connected and bipartite.
Consequently, among the four certified quartic orbital graphs, only a
connected bipartite member can possibly coincide with the graph under
consideration.

\begin{proposition}

Among the certified quartic orbital graphs on

\[
\Omega_{\mathrm{even}},
\]

exactly one is connected, bipartite, and triangle-free.

\end{proposition}

\begin{proof}

By the certified quartic orbital census of Theorem~\ref{thm:quartic-orbital-census}, exactly one
quartic orbital graph is connected.

The remaining three have two connected components of order fifteen.

The connected member is bipartite, whereas each disconnected member
contains triangles.

Therefore exactly one certified quartic orbital graph satisfies the
elementary graph-theoretic invariants established in Section~3.

\end{proof}

Accordingly, there is a unique certified candidate for identification with
the graph constructed in Section~3.

No graph isomorphism has yet been invoked.


\subsection{The coincidence theorem}

% STATUS: Conceptual deduction from certified computation
% SOURCE: graph isomorphism certificate

The graph constructed in Section~3 and the connected quartic orbital graph
identified in Section~5.4 arise from entirely different mathematical
constructions.

The first is obtained as the canonical bipartite double cover of the line
graph of the Petersen graph.

The second arises from the action of

\[
S_{5}
\]

on the homogeneous space

\[
S_{5}/V_{4,\mathrm{even}}.
\]

The following theorem identifies these constructions.

\begin{theorem}[Coincidence Theorem]
\label{thm:coincidence}

The connected triangle-free quartic orbital graph on

\[
S_{5}/V_{4,\mathrm{even}}
\]

is isomorphic to

\[
X=\KC(L(P)).
\]

\end{theorem}

\begin{proof}

The certified orbital census establishes the existence of exactly one
connected quartic orbital graph on

\[
S_{5}/V_{4,\mathrm{even}}.
\]

Section~3 constructs a connected bipartite quartic graph having the same
elementary graph-theoretic invariants.

A certified graph-isomorphism computation establishes an explicit
isomorphism between these two graphs.

Therefore the two constructions determine the same graph.

\end{proof}

Henceforth we make no distinction between these two presentations and
identify both constructions with the graph

\[
X.
\]
